% element: 10-s065:e04
% element: 10-s065:e05
\BeleskeID{10-s065}
Na granici dva režima važi jednakost raspoložive aktivne kolektorske struje i stvarne struje. Kontura u tom trenutku je
\[R_Bi_B(t_1)+V_{BE}=R_Ci_C(t_1)+V_{CES}=R_C\beta_Fi_B(t_1)+V_{CES}.\]
Rešavanjem se dobija bazna struja na slajdu. Izlaz u graničnom trenutku pišemo najpre kao
\[v_o(t_1)=V_{CC}-R_Bi_B(t_1)-V_{BE}-V_D,\]
a zatim zamenjujemo upravo dobijenu struju. Logaritamski izraz određuje vreme dostizanja te zadate izlazne vrednosti iz naznačene početne vrednosti i modelne asimptote. Pri primeni je presudno vezati vremensku konstantu i asimptotu za isti režim odziva.
